\documentclass[10pt,a4paper]{amsart}
\usepackage{amsmath,amssymb,amsthm,graphicx}
\usepackage{verbatim}
\newenvironment{reference}{\comment}{\endcomment}
%\newenvironment{reference}{}{}
\newenvironment{slogan}{\comment}{\endcomment}
\newenvironment{history}{\comment}{\endcomment}
% For commutative diagrams we use Xy-pic
\usepackage[all]{xy}

% We use 2cell for 2-commutative diagrams.
\xyoption{2cell}
% Theorem environments.
%
\theoremstyle{plain}
\newtheorem{theorem}[subsection]{Theorem}
\newtheorem{proposition}[subsection]{Proposition}
\newtheorem{lemma}[subsection]{Lemma}

\theoremstyle{definition}
\newtheorem{definition}[subsection]{Definition}
\newtheorem{example}[subsection]{Example}
\newtheorem{exercise}[subsection]{Exercise}
\newtheorem{situation}[subsection]{Situation}

\theoremstyle{remark}
\newtheorem{remark}[subsection]{Remark}
\newtheorem{remarks}[subsection]{Remarks}

\numberwithin{equation}{subsection}

% Macros
%
\def\lim{\mathop{\mathrm{lim}}\nolimits}
\def\colim{\mathop{\mathrm{colim}}\nolimits}
\def\Spec{\mathop{\mathrm{Spec}}}
\def\Hom{\mathop{\mathrm{Hom}}\nolimits}
\def\Ext{\mathop{\mathrm{Ext}}\nolimits}
\def\SheafHom{\mathop{\mathcal{H}\!\mathit{om}}\nolimits}
\def\SheafExt{\mathop{\mathcal{E}\!\mathit{xt}}\nolimits}
\def\Sch{\mathit{Sch}}
\def\Mor{\mathop{\mathrm{Mor}}\nolimits}
\def\Ob{\mathop{\mathrm{Ob}}\nolimits}
\def\Sh{\mathop{\mathit{Sh}}\nolimits}
\def\NL{\mathop{N\!L}\nolimits}
\def\CH{\mathop{\mathrm{CH}}\nolimits}
\def\proetale{{pro\text{-}\acute{e}tale}}
\def\etale{{\acute{e}tale}}
\def\QCoh{\mathit{QCoh}}
\def\Ker{\mathop{\mathrm{Ker}}}
\def\Im{\mathop{\mathrm{Im}}}
\def\Coker{\mathop{\mathrm{Coker}}}
\def\Coim{\mathop{\mathrm{Coim}}}

% Boxtimes
%
\DeclareMathSymbol{\boxtimes}{\mathbin}{AMSa}{"02}

%
% Macros for moduli stacks/spaces
%
\def\QCohstack{\mathcal{QC}\!\mathit{oh}}
\def\Cohstack{\mathcal{C}\!\mathit{oh}}
\def\Spacesstack{\mathcal{S}\!\mathit{paces}}
\def\Quotfunctor{\mathrm{Quot}}
\def\Hilbfunctor{\mathrm{Hilb}}
\def\Curvesstack{\mathcal{C}\!\mathit{urves}}
\def\Polarizedstack{\mathcal{P}\!\mathit{olarized}}
\def\Complexesstack{\mathcal{C}\!\mathit{omplexes}}
% \Pic is the operator that assigns to X its picard group, usage \Pic(X)
% \Picardstack_{X/B} denotes the Picard stack of X over B
% \Picardfunctor_{X/B} denotes the Picard functor of X over B
\def\Pic{\mathop{\mathrm{Pic}}\nolimits}
\def\Picardstack{\mathcal{P}\!\mathit{ic}}
\def\Picardfunctor{\mathrm{Pic}}
\def\Deformationcategory{\mathcal{D}\!\mathit{ef}}
\usepackage{cleveref}
\crefname{originallemma}{original lemma}{original lemmas}
\Crefname{originallemma}{Original lemma}{Original lemmas}
\makeatletter
\newlabel{lemma-construct-quasi-coherent-sheaves}{{lemma-construct-quasi-coherent-sheaves}{source}{}{sourcepointer-1}{}}
\newlabel{lemma-construct-quasi-coherent-sheaves@cref}{{[originallemma][0][]lemma-construct-quasi-coherent-sheaves}{[1][0][]source}{}{}{}}
\makeatother
\crefname{originalreference}{original reference}{original references}
\Crefname{originalreference}{Original reference}{Original references}
\makeatletter
\newlabel{sheaves-lemma-push-pull-composition-modules}{{sheaves-lemma-push-pull-composition-modules}{source}{}{sourcepointer-2}{}}
\newlabel{sheaves-lemma-push-pull-composition-modules@cref}{{[originalreference][0][]sheaves-lemma-push-pull-composition-modules}{[1][0][]source}{}{}{}}
\makeatother
\title{The pullback of the quasi-coherent module associated to a module is associated to the base changed module}
\author{The Stacks Project Authors}
\date{stacks-project:modules.tex:local-snapshot-20261012; GNU FDL 1.2 or later}
\begin{document}
\maketitle
\noindent\textit{Editorial note.} All delivered mathematical text is the author's own, copied from the pinned Stacks Project chapter source sources/stacks-modules/pinned-original-source.tex; no statement or proof step was rewritten, reordered or removed.\par\medskip
\noindent\textit{Editorial note.} Preamble provenance: the Stacks Project keeps one single shared preamble for every chapter; this card takes it directly from the staged sources/stacks-modules/preamble.tex as three verbatim spans (lines 8-12, 14-18 and 154-217), so no derived file was written and nothing was added or redefined. The wrapper class and the base packages are supplied by the assembler.\par\medskip
\noindent\textit{Editorial note.} Bibliography provenance: neither the counted statement nor its proof contains a citation (this lane rescans the whole span before writing the card), so the card carries no bibliography and no entry was reconstructed.\par\medskip
\noindent\textit{Editorial note.} Selection provenance: this card was opened under the file-based control contract (CONTROL.md), with the mechanical pre-card gate run on its statement and proof interval before the card was written.\par\medskip
\noindent\textit{Allowed standard prerequisites (editorial).}\begin{itemize}
\item The description of the quasi-coherent module associated to a module as the pullback along the structure morphism of the one point space. Invoked inside the counted proof at source lines 1448-1450 as a pointer to a lemma of this chapter. Reduces the statement to the compatibility of pullbacks with composition.
\item The commutativity of the square of ringed spaces formed by the structure morphisms and the given morphism. Written down in the counted proof (source lines 1452-1460), with a cross-chapter pointer for the composition of pullbacks. Identifies the two pullbacks of the module along the two routes of the square.
\end{itemize}\medskip
\ifdefined\hypertarget\hypertarget{sourcepointer-1}{}\fi
\noindent\textit{Source reference.} Original lemma lemma-construct-quasi-coherent-sheaves: sources/stacks-modules/pinned-original-source.tex, source lines 1425--1425. Navigation pointer only: the counted text refers to the Stacks item lemma-construct-quasi-coherent-sheaves, which belongs to another chapter of the project and is not part of this card. The pointer carries the target's own label as its number because this chapter's staged snapshot does not contain the other chapter's tag table; its content is not reproduced here.\par\medskip
\ifdefined\hypertarget\hypertarget{sourcepointer-2}{}\fi
\noindent\textit{Source reference.} Original reference sheaves-lemma-push-pull-composition-modules: sources/stacks-modules/pinned-original-source.tex, source lines 985--985. Navigation pointer only: the counted text refers to the Stacks item sheaves-lemma-push-pull-composition-modules, which belongs to another chapter of the project and is not part of this card. The pointer carries the target's own label as its number because this chapter's staged snapshot does not contain the other chapter's tag table; its content is not reproduced here.\par\medskip
\begin{lemma}
\label{lemma-restrict-quasi-coherent}
Let $(X, \mathcal{O}_X)$ be a ringed space.
Set $R = \Gamma(X, \mathcal{O}_X)$.
Let $M$ be an $R$-module.
Let $\mathcal{F}_M$ be the quasi-coherent sheaf of
$\mathcal{O}_X$-modules associated to $M$.
If $g : (Y, \mathcal{O}_Y) \to (X, \mathcal{O}_X)$
is a morphism of ringed spaces, then
$g^*\mathcal{F}_M$ is the sheaf associated
to the $\Gamma(Y, \mathcal{O}_Y)$-module
$\Gamma(Y, \mathcal{O}_Y) \otimes_R M$.
\end{lemma}

\begin{proof}
The assertion follows from the first description
of $\mathcal{F}_M$ in Lemma \ref{lemma-construct-quasi-coherent-sheaves}
as $\pi^*M$, and the following commutative diagram
of ringed spaces
$$
\xymatrix{
(Y, \mathcal{O}_Y) \ar[r]_-\pi \ar[d]_g &
(\{*\}, \Gamma(Y, \mathcal{O}_Y)) \ar[d]^{\text{induced by }g^\sharp} \\
(X, \mathcal{O}_X) \ar[r]^-\pi &
(\{*\}, \Gamma(X, \mathcal{O}_X))
}
$$
(Also use Sheaves, Lemma \ref{sheaves-lemma-push-pull-composition-modules}.)
\end{proof}

\end{document}
